% front/preface.tex —— 作者原序（Preface）
% 含：正文、引文、致谢、修订重印说明、第二版说明。原书末尾无署名。
% 注：每段必须写成一行（不得硬换行），否则 xeCJK 会把换行当成空格。
\cleardoublepage
\phantomsection
\chapter*{前言}
\addcontentsline{toc}{chapter}{前言}
\markboth{前言}{前言}

自 1946 年第一台现代计算机诞生以来, 数值线性代数 (又称矩阵计算) 一直是科学与工程计算的核心. 科学与工程中的大多数问题最终都归结为矩阵计算问题, 因此, 研究数值线性代数十分重要. 本书对矩阵计算作初步介绍, 同时也收录了近年来取得的一些新成果. 全书开篇, 我们首先在第 1 章给出数值线性代数的概貌.

第 2 章介绍求解一般线性方程组的基本直接法——高斯消去法. 高斯消去法通常用于求解规模中等, 没有特殊结构的稠密线性方程组, 其运算量为 $\mathcal{O}(n^{3})$, 其中 $n$ 为方程组的规模. 本章还讨论了主元技术.

第 3 章为讨论扰动与误差对数值解的影响, 引入向量范数与矩阵范数并研究其性质, 同时给出浮点运算与部分主元技术的误差分析.

第 4 章研究线性最小二乘问题, 重点讨论超定线性方程组 $A\mathbf{x}=\mathbf{b}$ ($A$ 的行数多于列数) 的最小二乘解的求法. 我们利用一些正交变换与 QR 分解来设计求解最小二乘问题的高效算法.

第 5 章研究求解 $A\mathbf{x}=\mathbf{b}$ 的经典迭代法. 迭代法与高斯消去法等直接法颇不相同. 当 $A$ 相当大时, 基于 $A$ 的 LU 分解的直接法在计算时间与计算机存储两方面都难以承受. 通常, 在大多数大型问题中矩阵是稀疏的; 但稀疏性可能在 LU 分解过程中丢失, 于是在 LU 分解结束时, 存储成为关键问题. 对这类问题, 可以使用一类称为迭代法的方法. 本章只讨论其中若干经典迭代法.

第 6 章介绍另一类迭代法, 即 Krylov 子空间方法. 我们只研究这类方法中的两个版本: 共轭梯度 (CG) 法与广义极小残量 (GMRES) 法. 1952 年提出的 CG 法是求解对称正定线性方程组最著名的迭代法之一; GMRES 方法则提出于 1986 年, 用于求解非对称线性方程组. 本章还简要讨论了预处理技术.

特征值问题在科学计算中格外引人关注. 第 7 章研究非对称特征值问题, 介绍幂法, 反幂法与 QR 方法等若干著名方法.

对称特征值问题具有良好的性质与丰富的数学理论, 是数值线性代数中最有趣的课题之一. 第 8 章研究这一课题, 讨论对称 QR 迭代法, Jacobi 法, 对分法以及分而治之技术.

第 9 章简要研究用边值方法求解带初值的常微分方程组 (ODEs) 的若干近期进展. 这类方法需要求解一个或多个非对称, 大型且稀疏的线性方程组, 因此我们用第 6 章的 GMRES 方法并配以适当的预条件子来求解它们. 主要结果之一是: 若对 $m$ 阶常微分方程组采用 $A_{\nu_1,\nu_2}$-稳定的边值方法, 则预处理后的矩阵可分解为 $I+L$, 其中 $I$ 为单位矩阵, $L$ 的秩至多为 $2m(\nu_1+\nu_2)$. 由此可知, 对预处理后的方程组施用 GMRES 方法, 至多经过 $2m(\nu_1+\nu_2)+1$ 次迭代即收敛. 文中还给出了在若干时滞微分方程上的应用.

第 10 章研究用于求解二维空间分数阶扩散方程的 Crank--Nicolson 加权与移位 Grünwald 差分 (CN-WSGD) 格式. 该格式需要求解一个或多个非对称, 大型且满的线性方程组. 我们使用第 6 章介绍的 GMRES 方法, 并结合重启策略与预处理技术来求解这些方程组. 主要结果之一是: 对预处理后的方程组施用重启 GMRES 方法时, 每个时间层的计算量为 $\mathcal{O}(N^{2}\log N)$, 其中 $N=N_{1}=N_{2}$ 为 $x$ 方向与 $y$ 方向的空间剖分数.

\vspace{0.8em}
\begin{quote}
「如果还有哪个数学分支之于数学科学的基础性, 像微积分与微分方程一样, 那便是数值线性代数. 」

\raggedleft ——\,L. Trefethen, D. Bau III
\end{quote}
\vspace{0.8em}

\noindent\textbf{致谢:} 我们要感谢香港中文大学数学系 Raymond H. F. Chan 教授始终如一的鼓励, 多年的友谊与经费支持; 感谢复旦大学数学科学学院 Z. H. Cao 教授的许多有益讨论与宝贵建议. 我们还要感谢朋友 Z. C. Shi 教授的鼓励与宝贵意见. 当然, 特别的谢意要献给作者生命中两个重要的机构——澳门大学与复旦大学, 它们为本书的写作提供了极好的学术氛围. 本书的大部分写作是在夜晚, 周末与假日完成的. 最后, 还要感谢我们的家人, 他们的关爱, 理解, 鼓励与支持是本书得以完成的必要条件. 谨向他们致以最诚挚的谢意!

\begingroup\raggedright
本书的出版部分得到澳门大学研究基金 (项目号 \mbox{RG024/01-02S/JXQ/FST}, \mbox{RG031/02-03S/JXQ/FST}, \mbox{RG064/03-04S/JXQ/FST}) 与国家自然科学基金项目 10471027 的资助; 上海市教育委员会与复旦大学亦提供了部分经费支持.
\par\endgroup

\medskip
\noindent\textbf{作者关于修订重印的说明:} 在本书的第二次印刷中, 我们订正了第一次印刷中的一些数学上与排印上的小错误. 谨向所有指出这些错误的读者致谢. 第 7 章作了若干补充与修订, 参考文献亦有更新, 书中还增加了一些习题. 本书的第二次印刷得到澳门大学研究基金 \mbox{RG081/04-05S/JXQ/FST} 的资助.

\medskip
\noindent\textbf{作者关于第二版的说明:} 自 2004 年本书问世至今已十余年, 我们收到了许多人的反馈. 同事, 学生与其他读者指出了书中一些数学上与排印上的小错误, 并对如何改进本书提出了建议. 我们尤其感谢上海海事大学数学系 Z. L. Xu 教授, 他提出了许多富有见地的意见与建设性的评论. 本版补入了定理~\ref{thm:6-7} 的证明, 并新增一章 (第 10 章). 第 10 章讨论数值分数阶扩散方程的近期进展. 参考文献有所扩充与更新. 本书的修订得到澳门大学研究基金 \mbox{MYRG098(Y3-L3)-FST13-JXQ}, 国家自然科学基金项目 11271084 以及复旦大学的资助.
